\chapter{2025年全国I卷}

\section{单选题}

\begin{problem}
\((1+5\i)\i\) 的虚部为（\quad）

\choices
    {\(-1\)}
    {\(0\)}
    {\(1\)}
    {\(6\)}
\end{problem}

\begin{problem}
设全集 \(U=\{x \mid x \text{ 是小于9的正整数}\}\)，集合 \(A=\{1,3,5\}\)，则 \(\complement_U A\) 中元素个数为（\quad）

\choices
    {\(0\)}
    {\(3\)}
    {\(5\)}
    {\(8\)}
\end{problem}

\begin{problem}
若双曲线 \(C\) 的虚轴长为实轴长的 \(\sqrt{7}\) 倍，则 \(C\) 的离心率为（\quad）

\choices
    {\(\sqrt{2}\)}
    {\(2\)}
    {\(\sqrt{7}\)}
    {\(2\sqrt{2}\)}
\end{problem}

\begin{problem}
若点 \((a,0)\)（\(a>0\)）是函数 \(y=2\tan\left(x-\frac{\pi}{3}\right)\) 的图象的一个对称中心，则 \(a\) 的最小值为（\quad）

\choices
    {\(\frac{\pi}{6}\)}
    {\(\frac{\pi}{3}\)}
    {\(\frac{\pi}{2}\)}
    {\(\frac{4\pi}{3}\)}
\end{problem}

\begin{problem}
设 \(f(x)\) 是定义在 \(\R\) 上且周期为2的偶函数，当 \(2 \le x \le 3\) 时，\(f(x)=5-2x\)，则 \(f\left(-\frac{3}{4}\right)=\)（\quad）

\choices
    {\(-\frac{1}{2}\)}
    {\(-\frac{1}{4}\)}
    {\(\frac{1}{4}\)}
    {\(\frac{1}{2}\)}
\end{problem}

\begin{problem}

帆船比赛中，运动员可借助风力计测定风速的大小与方向，测出的结果在航海学中称为视风风速，视风风速对应的向量是真风风速对应的向量与船行风速对应的向量之和，其中船行风速对应的向量与船速对应的向量大小相等，方向相反.
图1给出了部分风力等级，名称与风速大小的对应关系.
已知某帆船运动员在某时刻测得的视风风速对应的向量与船速对应的向量如图2所示（线段长度代表速度大小，单位：\(\mathrm{m/s}\)），则该时刻的真风为（\quad）

\begin{center}
\begin{tabular}{|c|c|c|}
\hline
级数 & 名称 & 风速大小（单位：\(\mathrm{m/s}\)） \\
\hline
2 & 轻风 & \(1.6\sim3.3\) \\
\hline
3 & 微风 & \(3.4\sim5.4\) \\
\hline
4 & 和风 & \(5.5\sim7.9\) \\
\hline
5 & 劲风 & \(8.0\sim10.7\) \\
\hline
\end{tabular}
\end{center}

\bitmapfigure[max width=6.0cm]{img/2025/national_paper_1/q6_fig2.png}

\choices
    {轻风}
    {微风}
    {和风}
    {劲风}
\end{problem}

\begin{problem}
若圆 \(x^2 + (y+2)^2 = r^2\)（\(r>0\)）上到直线 \(y = \sqrt{3}x + 2\) 的距离为1的点有且仅有2个，则 \(r\) 的取值范围是（\quad）

\choices
    {\((0,1)\)}
    {\((1,3)\)}
    {\((3,+\infty)\)}
    {\((0,+\infty)\)}
\end{problem}

\begin{problem}
若实数 \(x\)，\(y\)，\(z\) 满足 \(2 + \log_2 x = 3 + \log_3 y = 5 + \log_5 z\)，则 \(x\)，\(y\)，\(z\) 的大小关系不可能是（\quad）

\choices
    {\(x > y > z\)}
    {\(x > z > y\)}
    {\(y > x > z\)}
    {\(y > z > x\)}
\end{problem}

\section{多选题}

\begin{problem}
在正三棱柱 \(ABC - A_1B_1C_1\) 中，\(D\) 为 \(BC\) 中点，则（\quad）

\choices
    {\(AD \perp A_1C\)}
    {\(B_1C_1 \perp \text{平面} AA_1D\)}
    {\(AD \parallel A_1B_1\)}
    {\(CC_1 \parallel \text{平面} AA_1D\)}
\end{problem}

\begin{problem}
设抛物线 \(C: y^2 = 6x\) 的焦点为 \(F\)，过 \(F\) 的一条直线交 \(C\) 于 \(A\)，\(B\) 两点，过 \(A\) 作直线 \(l: x = -\frac{3}{2}\) 的垂线，垂足为 \(D\)，过 \(F\) 且与直线 \(AB\) 垂直的直线交 \(l\) 于 \(E\)，则（\quad）

\choices
    {\(|AD| = |AF|\)}
    {\(|AE| = |AB|\)}
    {\(|AB| \ge 6\)}
    {\(|AE| \cdot |BE| \ge 18\)}
\end{problem}

\begin{problem}
已知 \(\triangle ABC\) 的面积为 \(\frac{1}{4}\)，若 \(\cos 2A + \cos 2B + 2\sin C = 2\)，\(\cos A \cos B \sin C = \frac{1}{4}\)，则（\quad）

\choices
    {\(\sin C = \sin^2 A + \sin^2 B\)}
    {\(AB = \sqrt{2}\)}
    {\(\sin A + \sin B = \frac{\sqrt{6}}{2}\)}
    {\(AC^2 + BC^2 = 3\)}
\end{problem}

\section{填空题}

\begin{problem}
若直线 \(y = 2x + 5\) 是曲线 \(y = \e^x + x + a\) 的一条切线，则 \(a =\)\fillinblank{}.
\end{problem}

\begin{problem}
若一个等比数列的各项均为正数，且前4项和为4，前8项和为68，则该等比数列的公比为\fillinblank{}.
\end{problem}

\begin{problem}
有5个相同的球，分别标有数字 \(1\)，\(2\)，\(3\)，\(4\)，\(5\)，从中有放回地随机取3次，每次取1个球，记 \(X\) 为这5个球中至少被取出1次的球的个数，则数学期望 \(E(X) =\)\fillinblank{}.
\end{problem}

\section{解答题}

\begin{problem}
为研究某疾病与超声波检查结果的关系，从做过超声波检查的人群中随机调查了1000人，得到如下列联表：

\begin{center}
\begin{tabular}{|c|c|c|c|}
\hline
\multicolumn{1}{|c|}{组别} & \multicolumn{3}{c|}{超声波检查结果} \\
\cline{2-4}
 & 正常 & 不正常 & 合计 \\
\hline
患该疾病 & 20 & 180 & 200 \\
\hline
未患该疾病 & 780 & 20 & 800 \\
\hline
合计 & 800 & 200 & 1000 \\
\hline
\end{tabular}
\end{center}

\begin{enumerate}
\item 记超声波检查结果不正常者患该疾病的概率为 \(p\)，求 \(p\) 的估计值；
\item 根据小概率值 \(\alpha = 0.001\) 的独立性检验，分析超声波检查结果是否与患该疾病有关.
\end{enumerate}

附 \(\chi^2 = \frac{n(ad-bc)^2}{(a+b)(c+d)(a+c)(b+d)}\)

\begin{center}
\begin{tabular}{|c|c|c|c|}
\hline
\(P(\chi^2 \ge k)\) & 0.050 & 0.010 & 0.001 \\
\hline
\(k\) & 3.841 & 6.635 & 10.828 \\
\hline
\end{tabular}
\end{center}
\end{problem}

\begin{problem}
设数列 \(\{a_n\}\) 满足 \(a_1 = 3\)，\(\frac{a_{n+1}}{n} = \frac{a_n}{n+1} + \frac{1}{n(n+1)}\).

\begin{enumerate}
\item 证明：\(\{na_n\}\) 为等差数列；
\item 给定正整数 \(m\)，设函数 \(f(x) = a_1x + a_2x^2 + \cdots + a_mx^m\)，求 \(f'(-2)\).
\end{enumerate}
\end{problem}

\begin{problem}
如图所示的四棱锥 \(P - ABCD\) 中，\(PA \perp \text{平面} ABCD\)，\(BC \parallel AD\)，\(AB \perp AD\).

\begin{enumerate}
\item 证明：平面 \(PAB \perp \text{平面} PAD\)；
\item \(PA = AB = \sqrt{2}\)，\(BC = 2\)，\(AD = 1 + \sqrt{3}\)，点 \(P\)，\(B\)，\(C\)，\(D\) 均在球 \(O\) 的球面上.
\begin{enumerate}
\item 证明：\(O\) 在平面 \(ABCD\) 上；
\item 求直线 \(AC\) 与直线 \(PO\) 所成角的余弦值.
\end{enumerate}
\end{enumerate}

\bitmapfigure[max width=6.0cm]{img/2025/national_paper_1/q17_fig1.png}
\end{problem}

\begin{problem}
设椭圆 \(C: \frac{x^2}{a^2} + \frac{y^2}{b^2} = 1\)（\(a > b > 0\)） 的离心率为 \(\frac{2\sqrt{2}}{3}\)，下顶点为 \(A\)，右顶点为 \(B\)，\(|AB| = \sqrt{10}\).

\begin{enumerate}
\item 求椭圆 \(C\) 的标准方程；
\item 已知动点 \(P\) 不在 \(y\) 轴上，点 \(R\) 在射线 \(AP\) 上，且满足 \(|AP| \cdot |AR| = 3\).
\begin{enumerate}
\item 设 \(P(m, n)\)，求点 \(R\) 的坐标（用 \(m\)，\(n\) 表示）；
\item 设 \(O\) 为坐标原点，\(Q\) 是 \(C\) 上的动点，直线 \(OR\) 的斜率为直线 \(OP\) 的斜率的3倍，求 \(|PQ|\) 的最大值.
\end{enumerate}
\end{enumerate}
\end{problem}

\begin{problem}

\begin{enumerate}
\item 设函数 \(f(x) = 5\cos x - \cos 5x\)，求 \(f(x)\) 在 \(\left[0, \frac{\pi}{4}\right]\) 的最大值；
\item 给定 \(\theta \in (0,\pi)\) 和 \(a \in \R\)，证明：存在 \(y \in [a - \theta, a + \theta]\)，使得 \(\cos y \le \cos \theta\)；
\item 设 \(b \in \R\)，若存在 \(\varphi \in \R\) 使得 \(5\cos x - \cos (5x + \varphi) \le b\) 对 \(x \in \R\) 恒成立，求 \(b\) 的最小值.
\end{enumerate}
\end{problem}

